% Generated from validated Article Draft Result. Do not hand-edit; revise the structured draft instead.
\noindent\textbf{OpenAIは5月、モデルが離散幾何の長年の予想を反証したと報告した。同じ月には第三者評価のplaybookも公開された。単一の印象的成果を一般能力へ拡張せず、何をもって能力の証拠とするかを考える。}\autocite{src-1c93d6c1c725,src-8c84b8422093}

\subsection{5月の変化を一つのsystemとして読む}

standard benchmarkは比較可能性を与える一方、実際の研究課題で現れる能力をすべて覆うわけではない。5月20日にOpenAIが公表した離散幾何の事例は、その意味で強いattentionを集めるタイプの成果である。しかし、一次資料はOpenAI自身の説明であり、そこから「数学全般で人間を超えた」といった広い結論を導くことはできない。\autocite{src-1c93d6c1c725,src-8c84b8422093}


\subsection{An OpenAI model has disproved a central conjecture in discrete geometry}

OpenAIは5月20日、同社modelがunit-distance problemに関する長年の予想を反証したと報告した。公開日と「OpenAIがそう報告した」ことは確認できるが、数学的成果の意味やmodelの寄与範囲はfirst-party説明に依存する。本稿ではこの事例を、benchmark外のreal research taskで現れるcapabilityをどう検証するかという問題設定へつなげる。\autocite{src-1c93d6c1c725}


\subsection{A shared playbook for trustworthy third party evaluations}

5月29日、OpenAIはfrontier modelのthird-party evaluationについて、capability、safeguard、evaluation validityを扱うplaybookを公表した。評価者の独立性だけでなく、test conditionやmeasurement validityまで設計対象にする考え方は、印象的な単発成果を一般能力へ拡張しないための補助線になる。\autocite{src-8c84b8422093}


\subsection{Theme Synthesis}

重要なのは成果の派手さより、評価のchain of custodyである。誰が課題を設定し、model outputをどう検証し、第三者がどこまで再評価できるか。5月末のthird-party evaluation playbookは、capabilityだけでなくsafeguardやevaluation validityを評価対象に含める方向を示す。個別の成功例と一般的な能力評価を分離することが、frontier modelを読むうえでますます重要になっている。\autocite{src-1c93d6c1c725,src-8c84b8422093}


\begin{claimboundary}
Claim Boundary: first-party sourceは公開日とsource ownerが説明したscopeを確認する一次資料だが、capability・benchmark・safety・performanceの優位性まで独立検証したものではない。\autocite{src-1c93d6c1c725,src-8c84b8422093}
\end{claimboundary}
